\documentclass[10pt]{article}
\usepackage[T1]{fontenc}
\usepackage[utf8]{inputenc}
\usepackage[french]{babel}
\usepackage{amsmath,amssymb,mathrsfs,mathpazo}
\usepackage{geometry}
\geometry{a4paper,margin=2cm}
\usepackage{graphicx,array,tabularx,tabularray,colortbl,multirow,multicol,booktabs,enumitem,pifont,ulem,fancybox,diagbox,cancel}
\usepackage{pstricks,pstricks-add,pst-plot,pst-tree,pst-node,pst-text,pst-func,pst-math,pst-bspline,pst-3dplot,pst-eucl,pst-solides3d}
\let\pscalculate\undefined
\usepackage{pst-bezier}
\usepackage{tikz,pgfplots,tkz-tab}
\usetikzlibrary{arrows,arrows.meta,patterns,calc,positioning,trees,decorations.markings,shapes.geometric,shapes.misc,shapes.arrows,intersections,decorations.pathmorphing,decorations.pathreplacing,backgrounds,shadows,babel}
\pgfplotsset{compat=1.18}
\usepgfplotslibrary{fillbetween}
\usepackage[np]{numprint}
\usepackage{listings,makecell,pas-tableur,twoopt}
\let\barre\undefined
\usepackage[pstricks]{tablvar}
\definecolor{darkgreen}{RGB}{0,100,0}
\providecommand{\red}{\color{red}}
\providecommand{\blue}{\color{blue}}
\providecommand{\green}{\color{green}}
\providecommand{\black}{\color{black}}
\providecommand{\white}{\color{white}}
\providecommand{\gray}{\color{gray}}
\providecommand{\cyan}{\color{cyan}}
\providecommand{\magenta}{\color{magenta}}
\providecommand{\yellow}{\color{yellow}}
\providecommand{\violet}{\color{violet}}
\providecommand{\orange}{\color{orange}}
\providecommand{\euro}{\texteuro}
\providecommand{\EUR}{\texteuro}
\providecommand{\ieme}{\textsuperscript{e}}
\providecommand{\degres}{\ensuremath{^\circ}}
\providecommand{\degree}{\ensuremath{^\circ}}
\providecommand{\vv}[1]{\overrightarrow{#1}}
\providecommand{\Coord}[1]{#1}
\providecommand{\decosix}{}
\providecommand{\starredbullet}{\ensuremath{\bullet}}
\pagestyle{empty}
\setlength{\parindent}{0pt}
\usepackage{adjustbox}\DeclareUnicodeCharacter{2236}{:}\DeclareUnicodeCharacter{2190}{\ensuremath{\leftarrow}}
\providecommand{\R}{\mathbb{R}}
\providecommand{\N}{\mathbb{N}}
\providecommand{\Z}{\mathbb{Z}}
\providecommand{\Q}{\mathbb{Q}}
\providecommand{\D}{\mathbb{D}}
\providecommand{\C}{\mathbb{C}}
\providecommand{\e}{{\mathrm{e}}}
\providecommand{\vect}[1]{\overrightarrow{#1}}
\providecommand{\vectt}[1]{\overrightarrow{\mathrm{#1}}}
\providecommand{\vvt}[1]{\overrightarrow{\mathrm{#1}}}
\providecommand{\np}[2][]{{#2~#1}}
\providecommand{\vv}[1]{{\overrightarrow{#1}}}
\renewcommand{\d}{{\mathrm{d}}}
\providecommand{\degres}{{^\circ}}
\providecommand{\degre}{{^\circ}}
\providecommand{\blue}{\color{blue}}
\providecommand{\red}{\color{red}}
\providecommand{\enskip}{\quad}
\providecommand{\llbracket}{\lbrack\!\lbrack}
\providecommand{\rrbracket}{\rbrack\!\rbrack}
\providecommand{\euro}{{\text{€}}}
\providecommand{\no}{{\text{n°}}}
\providecommand{\interval}[2]{{[#1;#2]}}
\providecommand{\intervFO}[2]{{[#1;#2[}}
\providecommand{\dotfill}{\ldots}

\providecommand{\ds}{\displaystyle}
\providecommand{\textding}[1]{#1}
\providecommand{\Oij}{\(\mathrm{O};\vec{\imath},\vec{\jmath}\)}
\providecommand{\Oijk}{\(\mathrm{O};\vec{\imath},\vec{\jmath},\vec{k}\)}
\providecommand{\Ouv}{\(\mathrm{O};\vec{u},\vec{v}\)}
\providecommand{\intervFF}[2]{\ensuremath{[#1;#2]}}
\renewcommand{\euro}{\ifmmode\text{\texteuro}\else\texteuro\fi}
\DeclareUnicodeCharacter{00B0}{\ifmmode^{\circ}\else\textdegree\fi}
\DeclareUnicodeCharacter{20AC}{\ifmmode\text{\texteuro}\else\texteuro\fi}
% Énoncés : Éducation nationale. Transcription et corrigés : APMEP, Denis Vergès et contributeurs. Découpage et classement : CoopMaths, Sébastien Lozano, Éric Elter et contributeurs. Adaptation de la présentation : Maths Explorer.
% CC BY-SA : https://coopmaths.fr/www/about/
% https://coopmaths.fr/alea/static/bac/2021/tex/bac_2021_03_sujet2_3.tex
\begin{document}
\noindent\textit{APMEP / CoopMaths. CC BY-SA. Présentation adaptée par Maths Explorer.}\par\medskip
\medskip

Dans l'espace rapporté à un repère orthonormé \Oijk, on considère les points:

\text{A} de coordonnées (2~;~0~;~0), B de coordonnées (0~;~3~;~0) et C de coordonnées
(0~;~0~;~1).

\psset{unit=1cm}
\begin{center}
\begin{pspicture}(-0.5,-0.5)(12,6)
\pspolygon[fillstyle=solid,fillcolor=gray!20,linestyle=dashed](0,0)(2.4,4.5)(10.4,2)

\psline(0,0)(8,0)(10.4,2)(10.4,4.5)(8,2.5)(8,0)
\psline(8,2.5)(0,2.5)(0,0)
\psline(0,2.5)(2.4,4.5)(10.4,4.5)
\psline[linestyle=dashed](0,0)(2.4,2)(2.4,4.5)
\psline[linestyle=dashed](2.4,2)(10.4,2)
\uput[dl](0,0){A}\uput[u](2.4,4.5){C}\uput[r](10.4,2){B}\uput[d](2.5,2){O}
\end{pspicture}
\end{center}

\medskip

L'objectif de cet exercice est de calculer l'aire du triangle ABC.

\medskip

\begin{enumerate}
\item 
	\begin{enumerate}
		\item Montrer que le vecteur $\vect{n}\begin{pmatrix}3\\2\\6\end{pmatrix}$ est normal au plan (ABC).
		\item En déduire qu'une équation cartésienne du plan (ABC) est : $3x + 2y + 6z - 6 = 0$.
	\end{enumerate}
\item  On note $d$ la droite passant par O et orthogonale au plan (ABC). 
	\begin{enumerate}
		\item Déterminer une représentation paramétrique de la droite $d$.
		\item Montrer que la droite $d$ coupe le plan (ABC) au point H de coordonnées $\left(\frac{18}{49}~;~\frac{12}{49}~;~\frac{36}{49}\right)$.
		\item Calculer la distance OH.
	\end{enumerate}
\item  On rappelle que le volume d'une pyramide est donné par: $V = \dfrac{1}{3}\mathcal{B}h$, où $\mathcal{B}$ est l'aire d'une
base et $h$ est la hauteur de la pyramide correspondant à cette base.

En calculant de deux façons différentes le volume de la pyramide OABC, déterminer l'aire du triangle ABC.
\end{enumerate}

\bigskip

\end{document}
